%! TEX root = ../m1-20-21.tex
\chapter{Serii de numere reale pozitive}
\index{serii!de numere pozitive}

Intuitiv, o serie poate fi gîndită ca o sumă infinită, dată de o
regulă a unui termen general. De exemplu, seria:
\[
    \sum_{n \geq 0} \dfrac{3n}{5n^2 + 2n - 1}
\]
are termenul general de forma $ x_n = \dfrac{3n}{5n^2 + 2n - 1} $
și putem rescrie seria mai simplu $ \sum x_n $, presupunînd implicit
că indicele $ n $ ia cea mai mică valoare permisă și merge pînă la $ \infty $.

Natura seriilor este fundamental diferită de cea a șirurilor prin faptul că
seriile \emph{acumulează}. De exemplu, să considerăm șirul constant $ a_n = 1, \forall n $.
Atunci, evident, $ \ds\lim_{n \to \infty} a_n = 1 $. Pe de altă parte, dacă luăm
seria de termen general $ a_n $, adică $ \sum a_n $, observăm că aceasta are suma $ \infty $,
deci este divergentă.

În continuare, vom studia criterii prin care putem decide dacă o serie este
sau nu convergentă. Dar înainte de aceasta, vom folosi foarte des două serii
particulare, pe care le detaliem în continuare.

%%%%%%%%%%%%%%%%%%%%%%%%%%%%%%%%%%%%%%%%%%%%%%%%%%%%%%%%%%%%%%%%%%%%%%
\section{Seria geometrică}
\index{serii!geometrice}

Pornim de la progresiile geometrice studiate în liceu. Fie $ (b_n) $ o progresie
geometrică, cu primul termen $ b_1 $ și cu rația $ q $. Deci termenul general
are formula $ b_n = b_1 q^{n-1} $. Atunci suma primilor $ n $ termeni ai progresiei
se poate calcula cu formula:
\[
    S_n = \sum_{k = 1}^n b_k = b_1 \cdot \dfrac{q^n - 1}{q - 1}.
\]

Dacă, însă, în această sumă considerăm \qq{toți} termenii progresiei, obținem
\emph{seria} geometrică, anume $ \ds\sum_{n \geq 1} b_n $.

Suma acestei serii coincide cu $ \ds\lim_{n \to \infty} S_n $ și se poate
observa cu ușurință că seria geometrică este convergentă dacă și numai dacă
$ |q| < 1 $. Mai mult, în caz de convergență, suma seriei se poate calcula
imediat ca fiind $ b_1 \cdot \dfrac{1}{1 - q} $.

%%%%%%%%%%%%%%%%%%%%%%%%%%%%%%%%%%%%%%%%%%%%%%%%%%%%%%%%%%%%%%%%%%%%%%
\section{Seria armonică}
\index{serii!armonice}

Această serie se mai numește \emph{funcția zeta a lui Riemann} și se
definește astfel:
\[
    \zeta(s) = \sum_{n \geq 1} \dfrac{1}{n^s}, \quad s \in \QQ.
\]

Remarcăm cîteva cazuri particulare:
\begin{align*}
    \zeta(0) &= \sum 1 = \infty \\
    \zeta(1) &= \sum \dfrac{1}{n} = \infty \\
    \zeta(2) &= \sum \dfrac{1}{n^2} = \dfrac{\pi^2}{6} \\
    \zeta(-1) &= \sum n = \infty \\
    \zeta(-2) &= \sum n^2 = \infty.
\end{align*}

Rezultatul general este:
\begin{theorem}\label{thm:zeta-conv}
    Seria armonică $ \zeta(s) $ este convergentă dacă și numai dacă
    $ s > 1 $.
\end{theorem}

%%%%%%%%%%%%%%%%%%%%%%%%%%%%%%%%%%%%%%%%%%%%%%%%%%%%%%%%%%%%%%%%%%%%%%
\section{Criterii de convergență}

Ne păstrăm în continuare în contextul seriilor cu termeni reali și pozitivi,
pe care le scriem în general $ \sum x_n $.

Convergența poate fi decisă ușor folosind criteriile de mai jos.

\vspace{1cm}

\index{criteriul!necesar}
\textbf{Criteriul necesar:} Dacă $ \ds\lim_{n \to \infty} x_n \neq 0 $, atunci
seria $ \sum x_n $ este divergentă.

\begin{remark}\label{rk:nec}
    Să remarcăm că, așa cum îi spune și numele, criteriul de mai sus dă doar
    condiții necesare, nu și suficiente pentru convergență! De exemplu, pentru
    $ \zeta(1) $, termenul general tinde către 0, dar seria este divergentă.
\end{remark}

\vspace{1cm}

\index{criteriul!de comparație!termen cu termen}
\textbf{Criteriul de comparație termen cu termen:} Fie $ \sum y_n $ o altă serie
de numere reale și pozitive.
\begin{itemize}
    \item Dacă $ x_n \leq y_n, \forall n $, iar seria $ \sum y_n $ este convergentă,
        atunci și seria $ \sum x_n $ este convergentă;
    \item Dacă $ x_n \geq y_n, \forall n $, iar seria $ \sum y_n $ este divergentă,
        atunci și seria $ \sum x_n $ este divergentă.
\end{itemize}

Acest criteriu seamănă foarte mult cu criteriul de comparație de la șiruri.
Astfel, avem că un șir mai mare (termen cu termen) decît un șir divergent este
divergent, iar un șir mai mic (termen cu termen) decît un șir convergent este
convergent. Celelalte cazuri sînt nedecise.

De asemenea, mai remarcăm că, în studiul seriei $ \sum x_n $ apare seria $ \sum y_n $,
care trebuie aleasă convenabil astfel încît să aibă loc condițiile criteriului.
În practică, cel mai des vom alege această nouă serie ca fiind o serie geometrică
sau una armonică, cu rația, respectiv exponentul alese convenabil.

\vspace{1cm}

\index{criteriul!de comparație!la limită}
\textbf{Criteriul de comparație la limită:} Fie $ \sum y_n $ o altă serie de numere
reale și pozitive, astfel încît $ \ds\lim_{n \to \infty} \dfrac{x_n}{y_n} \in (0, \infty) $.
Atunci cele două serii au aceeași natură, adică $ \sum x_n $ este convergentă
dacă și numai dacă $ \ds\sum y_n $ este convergentă.

\vspace{1cm}

\index{criteriul!raportului}
\textbf{Criteriul raportului:} Fie $ \ell = \ds\lim_{n \to \infty} \dfrac{x_{n+1}}{x_n} $.
\begin{itemize}
    \item Dacă $ \ell > 1 $, atunci seria $ \sum x_n $ este divergentă;
    \item Dacă $ \ell < 1 $, atunci seria $ \sum x_n $ este convergentă;
    \item Dacă $ \ell = 1 $, atunci criteriul nu decide.
\end{itemize}

\vspace{1cm}

\index{criteriul!radical}
\textbf{Criteriul radical:} Fie $ \ell = \ds\lim_{n \to \infty} \sqrt[n]{x_n} $.
\begin{itemize}
    \item Dacă $ \ell > 1 $, atunci seria $ \sum x_n $ este divergentă;
    \item Dacă $ \ell < 1 $, atunci seria $ \sum x_n $ este convergentă;
    \item Dacă $ \ell = 1 $, atunci criteriul nu decide.
\end{itemize}

%%%%%%%%%%%%%%%%%%%%%%%%%%%%%%%%%%%%%%%%%%%%%%%%%%%%%%%%%%%%%%%%%%%%%%
\newpage
\section{Exerciții}

1. Studiați convergența seriilor $ \sum x_n $ în cazurile de mai jos:
\begin{enumerate}[(a)]
    \item $ x_n = \left( \dfrac{3n}{3n + 1} \right)^n $; \hfill(D, necesar)
    \item $ x_n = \dfrac{1}{n!} $; \hfill(C, raport)
    \item $ x_n = \dfrac{1}{n\sqrt{n + 1}} $; \hfill(C, comparație la limită cu $\zeta(3/2)$)
    \item $ x_n = \arcsin \dfrac{n+1}{2n + 3} $; \hfill(D, necesar)
    \item $ x_n = \left( 1 - \dfrac{1}{n} \right)^n $; \hfill(D, necesar)
    \item $ x_n = \dfrac{n!}{n^{2n}} $; \hfill(C, raport)
    \item $ x_n = \left( \dfrac{n + 1}{3n + 1}\right)^n $; \hfill(C, radical)
    \item $ x_n = \left( 1 - \dfrac{1}{n} \right)^{n^2} $; \hfill(C, radical)
    \item $ x_n = \dfrac{1}{7^n + 3^n} $; \hfill(C, comparație cu geometrică)
    \item $ x_n = \dfrac{2 + \sin n}{n^2} $; \hfill(C, comparație cu armonică)
    \item $ x_n = \dfrac{\sin^2 n}{n^2 + 1} $; \hfill(C, comparație cu armonică)
    \item $ x_n = \sqrt{n^4 + 2n + 1} - n^2 $; \hfill(D, comparație cu armonică)
\end{enumerate}

2*. Studiați convergența \emph{șirurilor} cu termenul general $ x_n $ din
exercițiul anterior și comparați cu comportamentul seriilor.

%%% Local Variables:
%%% mode: latex
%%% TeX-master: "../m1-20-21"
%%% End:
